\IfFileExists{stacks-project.cls}{%
\documentclass{stacks-project}
}{%
\documentclass{amsart}
}

% For dealing with references we use the comment environment
\usepackage{verbatim}
\newenvironment{reference}{\comment}{\endcomment}
%\newenvironment{reference}{}{}
\newenvironment{slogan}{\comment}{\endcomment}
\newenvironment{history}{\comment}{\endcomment}

% For commutative diagrams we use Xy-pic
\usepackage[all]{xy}

% We use 2cell for 2-commutative diagrams.
\xyoption{2cell}
\UseAllTwocells

% We use multicol for the list of chapters between chapters
\usepackage{multicol}

% This is generally recommended for better output
\usepackage{lmodern}
\usepackage[T1]{fontenc}

% For cross-file-references
\usepackage{xr-hyper}

% Package for hypertext links:
\usepackage{hyperref}

% For any local file, say "hello.tex" you want to link to please
% use \externaldocument[hello-]{hello}
\externaldocument[introduction-]{introduction}
\externaldocument[conventions-]{conventions}
\externaldocument[sets-]{sets}
\externaldocument[categories-]{categories}
\externaldocument[topology-]{topology}
\externaldocument[sheaves-]{sheaves}
\externaldocument[sites-]{sites}
\externaldocument[stacks-]{stacks}
\externaldocument[fields-]{fields}
\externaldocument[algebra-]{algebra}
\externaldocument[brauer-]{brauer}
\externaldocument[homology-]{homology}
\externaldocument[derived-]{derived}
\externaldocument[simplicial-]{simplicial}
\externaldocument[more-algebra-]{more-algebra}
\externaldocument[smoothing-]{smoothing}
\externaldocument[modules-]{modules}
\externaldocument[sites-modules-]{sites-modules}
\externaldocument[injectives-]{injectives}
\externaldocument[cohomology-]{cohomology}
\externaldocument[sites-cohomology-]{sites-cohomology}
\externaldocument[dga-]{dga}
\externaldocument[dpa-]{dpa}
\externaldocument[sdga-]{sdga}
\externaldocument[hypercovering-]{hypercovering}
\externaldocument[schemes-]{schemes}
\externaldocument[constructions-]{constructions}
\externaldocument[properties-]{properties}
\externaldocument[morphisms-]{morphisms}
\externaldocument[coherent-]{coherent}
\externaldocument[divisors-]{divisors}
\externaldocument[limits-]{limits}
\externaldocument[varieties-]{varieties}
\externaldocument[topologies-]{topologies}
\externaldocument[descent-]{descent}
\externaldocument[perfect-]{perfect}
\externaldocument[more-morphisms-]{more-morphisms}
\externaldocument[flat-]{flat}
\externaldocument[groupoids-]{groupoids}
\externaldocument[more-groupoids-]{more-groupoids}
\externaldocument[etale-]{etale}
\externaldocument[chow-]{chow}
\externaldocument[intersection-]{intersection}
\externaldocument[pic-]{pic}
\externaldocument[weil-]{weil}
\externaldocument[adequate-]{adequate}
\externaldocument[dualizing-]{dualizing}
\externaldocument[duality-]{duality}
\externaldocument[discriminant-]{discriminant}
\externaldocument[derham-]{derham}
\externaldocument[local-cohomology-]{local-cohomology}
\externaldocument[algebraization-]{algebraization}
\externaldocument[curves-]{curves}
\externaldocument[resolve-]{resolve}
\externaldocument[models-]{models}
\externaldocument[functors-]{functors}
\externaldocument[equiv-]{equiv}
\externaldocument[pione-]{pione}
\externaldocument[etale-cohomology-]{etale-cohomology}
\externaldocument[proetale-]{proetale}
\externaldocument[relative-cycles-]{relative-cycles}
\externaldocument[more-etale-]{more-etale}
\externaldocument[trace-]{trace}
\externaldocument[crystalline-]{crystalline}
\externaldocument[spaces-]{spaces}
\externaldocument[spaces-properties-]{spaces-properties}
\externaldocument[spaces-morphisms-]{spaces-morphisms}
\externaldocument[decent-spaces-]{decent-spaces}
\externaldocument[spaces-cohomology-]{spaces-cohomology}
\externaldocument[spaces-limits-]{spaces-limits}
\externaldocument[spaces-divisors-]{spaces-divisors}
\externaldocument[spaces-over-fields-]{spaces-over-fields}
\externaldocument[spaces-topologies-]{spaces-topologies}
\externaldocument[spaces-descent-]{spaces-descent}
\externaldocument[spaces-perfect-]{spaces-perfect}
\externaldocument[spaces-more-morphisms-]{spaces-more-morphisms}
\externaldocument[spaces-flat-]{spaces-flat}
\externaldocument[spaces-groupoids-]{spaces-groupoids}
\externaldocument[spaces-more-groupoids-]{spaces-more-groupoids}
\externaldocument[bootstrap-]{bootstrap}
\externaldocument[spaces-pushouts-]{spaces-pushouts}
\externaldocument[spaces-chow-]{spaces-chow}
\externaldocument[groupoids-quotients-]{groupoids-quotients}
\externaldocument[spaces-more-cohomology-]{spaces-more-cohomology}
\externaldocument[spaces-simplicial-]{spaces-simplicial}
\externaldocument[spaces-duality-]{spaces-duality}
\externaldocument[formal-spaces-]{formal-spaces}
\externaldocument[restricted-]{restricted}
\externaldocument[spaces-resolve-]{spaces-resolve}
\externaldocument[formal-defos-]{formal-defos}
\externaldocument[defos-]{defos}
\externaldocument[cotangent-]{cotangent}
\externaldocument[examples-defos-]{examples-defos}
\externaldocument[algebraic-]{algebraic}
\externaldocument[examples-stacks-]{examples-stacks}
\externaldocument[stacks-sheaves-]{stacks-sheaves}
\externaldocument[criteria-]{criteria}
\externaldocument[artin-]{artin}
\externaldocument[quot-]{quot}
\externaldocument[stacks-properties-]{stacks-properties}
\externaldocument[stacks-morphisms-]{stacks-morphisms}
\externaldocument[stacks-limits-]{stacks-limits}
\externaldocument[stacks-cohomology-]{stacks-cohomology}
\externaldocument[stacks-perfect-]{stacks-perfect}
\externaldocument[stacks-introduction-]{stacks-introduction}
\externaldocument[stacks-more-morphisms-]{stacks-more-morphisms}
\externaldocument[stacks-geometry-]{stacks-geometry}
\externaldocument[moduli-]{moduli}
\externaldocument[moduli-curves-]{moduli-curves}
\externaldocument[examples-]{examples}
\externaldocument[exercises-]{exercises}
\externaldocument[guide-]{guide}
\externaldocument[desirables-]{desirables}
\externaldocument[coding-]{coding}
\externaldocument[obsolete-]{obsolete}
\externaldocument[fdl-]{fdl}
\externaldocument[index-]{index}

% Theorem environments.
%
\theoremstyle{plain}
\newtheorem{theorem}[subsection]{Theorem}
\newtheorem{proposition}[subsection]{Proposition}
\newtheorem{lemma}[subsection]{Lemma}

\theoremstyle{definition}
\newtheorem{definition}[subsection]{Definition}
\newtheorem{example}[subsection]{Example}
\newtheorem{exercise}[subsection]{Exercise}
\newtheorem{situation}[subsection]{Situation}

\theoremstyle{remark}
\newtheorem{remark}[subsection]{Remark}
\newtheorem{remarks}[subsection]{Remarks}

\numberwithin{equation}{subsection}

% Macros
%
\def\lim{\mathop{\mathrm{lim}}\nolimits}
\def\colim{\mathop{\mathrm{colim}}\nolimits}
\def\Spec{\mathop{\mathrm{Spec}}}
\def\Hom{\mathop{\mathrm{Hom}}\nolimits}
\def\Ext{\mathop{\mathrm{Ext}}\nolimits}
\def\SheafHom{\mathop{\mathcal{H}\!\mathit{om}}\nolimits}
\def\SheafExt{\mathop{\mathcal{E}\!\mathit{xt}}\nolimits}
\def\Sch{\mathit{Sch}}
\def\Mor{\mathop{\mathrm{Mor}}\nolimits}
\def\Ob{\mathop{\mathrm{Ob}}\nolimits}
\def\Sh{\mathop{\mathit{Sh}}\nolimits}
\def\NL{\mathop{N\!L}\nolimits}
\def\CH{\mathop{\mathrm{CH}}\nolimits}
\def\proetale{{pro\text{-}\acute{e}tale}}
\def\etale{{\acute{e}tale}}
\def\QCoh{\mathit{QCoh}}
\def\Ker{\mathop{\mathrm{Ker}}}
\def\Im{\mathop{\mathrm{Im}}}
\def\Coker{\mathop{\mathrm{Coker}}}
\def\Coim{\mathop{\mathrm{Coim}}}

% Boxtimes
%
\DeclareMathSymbol{\boxtimes}{\mathbin}{AMSa}{"02}

%
% Macros for moduli stacks/spaces
%
\def\QCohstack{\mathcal{QC}\!\mathit{oh}}
\def\Cohstack{\mathcal{C}\!\mathit{oh}}
\def\Spacesstack{\mathcal{S}\!\mathit{paces}}
\def\Quotfunctor{\mathrm{Quot}}
\def\Hilbfunctor{\mathrm{Hilb}}
\def\Curvesstack{\mathcal{C}\!\mathit{urves}}
\def\Polarizedstack{\mathcal{P}\!\mathit{olarized}}
\def\Complexesstack{\mathcal{C}\!\mathit{omplexes}}
% \Pic is the operator that assigns to X its picard group, usage \Pic(X)
% \Picardstack_{X/B} denotes the Picard stack of X over B
% \Picardfunctor_{X/B} denotes the Picard functor of X over B
\def\Pic{\mathop{\mathrm{Pic}}\nolimits}
\def\Picardstack{\mathcal{P}\!\mathit{ic}}
\def\Picardfunctor{\mathrm{Pic}}
\def\Deformationcategory{\mathcal{D}\!\mathit{ef}}
\usepackage{cleveref}
\crefname{originalreference}{original reference}{original references}
\Crefname{originalreference}{Original reference}{Original references}
\makeatletter
\newlabel{algebraic-definition-algebraic-stack}{{026O}{source}{}{sourcepointer-1}{}}
\newlabel{algebraic-definition-algebraic-stack@cref}{{[originalreference][0][]026O}{[1][0][]source}{}{}{}}
\makeatother
\crefname{originalreference}{original reference}{original references}
\Crefname{originalreference}{Original reference}{Original references}
\makeatletter
\newlabel{algebraic-lemma-map-space-into-stack}{{04T4}{source}{}{sourcepointer-2}{}}
\newlabel{algebraic-lemma-map-space-into-stack@cref}{{[originalreference][0][]04T4}{[1][0][]source}{}{}{}}
\makeatother
\crefname{originalreference}{original reference}{original references}
\Crefname{originalreference}{Original reference}{Original references}
\makeatletter
\newlabel{algebraic-lemma-stack-presentation}{{04T5}{source}{}{sourcepointer-3}{}}
\newlabel{algebraic-lemma-stack-presentation@cref}{{[originalreference][0][]04T5}{[1][0][]source}{}{}{}}
\makeatother
\crefname{originalreference}{original reference}{original references}
\Crefname{originalreference}{Original reference}{Original references}
\makeatletter
\newlabel{algebraic-theorem-smooth-groupoid-gives-algebraic-stack}{{04TK}{source}{}{sourcepointer-4}{}}
\newlabel{algebraic-theorem-smooth-groupoid-gives-algebraic-stack@cref}{{[originalreference][0][]04TK}{[1][0][]source}{}{}{}}
\makeatother
\crefname{originalreference}{original reference}{original references}
\Crefname{originalreference}{Original reference}{Original references}
\makeatletter
\newlabel{bootstrap-lemma-spaces-etale}{{076M}{source}{}{sourcepointer-5}{}}
\newlabel{bootstrap-lemma-spaces-etale@cref}{{[originalreference][0][]076M}{[1][0][]source}{}{}{}}
\makeatother
\crefname{originalreference}{original reference}{original references}
\Crefname{originalreference}{Original reference}{Original references}
\makeatletter
\newlabel{more-morphisms-lemma-etale-dominates-smooth}{{055V}{source}{}{sourcepointer-6}{}}
\newlabel{more-morphisms-lemma-etale-dominates-smooth@cref}{{[originalreference][0][]055V}{[1][0][]source}{}{}{}}
\makeatother
\crefname{originalreference}{original reference}{original references}
\Crefname{originalreference}{Original reference}{Original references}
\makeatletter
\newlabel{spaces-definition-algebraic-space}{{025Y}{source}{}{sourcepointer-7}{}}
\newlabel{spaces-definition-algebraic-space@cref}{{[originalreference][0][]025Y}{[1][0][]source}{}{}{}}
\makeatother
\crefname{originalreference}{original reference}{original references}
\Crefname{originalreference}{Original reference}{Original references}
\makeatletter
\newlabel{spaces-groupoids-section-explicit-quotient-stacks}{{04MB}{source}{}{sourcepointer-8}{}}
\newlabel{spaces-groupoids-section-explicit-quotient-stacks@cref}{{[originalreference][0][]04MB}{[1][0][]source}{}{}{}}
\makeatother
\crefname{originalreference}{original reference}{original references}
\Crefname{originalreference}{Original reference}{Original references}
\makeatletter
\newlabel{stacks-lemma-characterize-essentially-surjective-when-ff}{{046N}{source}{}{sourcepointer-9}{}}
\newlabel{stacks-lemma-characterize-essentially-surjective-when-ff@cref}{{[originalreference][0][]046N}{[1][0][]source}{}{}{}}
\makeatother
\crefname{originalreference}{original reference}{original references}
\Crefname{originalreference}{Original reference}{Original references}
\makeatletter
\newlabel{stacks-sheaves-lemma-surjective-flat-locally-finite-presentation}{{06XH}{source}{}{sourcepointer-10}{}}
\newlabel{stacks-sheaves-lemma-surjective-flat-locally-finite-presentation@cref}{{[originalreference][0][]06XH}{[1][0][]source}{}{}{}}
\makeatother
\crefname{originalreference}{original reference}{original references}
\Crefname{originalreference}{Original reference}{Original references}
\makeatletter
\newlabel{stacks-sheaves-section-sheaves}{{06TN}{source}{}{sourcepointer-11}{}}
\newlabel{stacks-sheaves-section-sheaves@cref}{{[originalreference][0][]06TN}{[1][0][]source}{}{}{}}
\makeatother
\crefname{originalreference}{original reference}{original references}
\Crefname{originalreference}{Original reference}{Original references}
\makeatletter
\newlabel{topologies-remark-choice-sites}{{03FF}{source}{}{sourcepointer-12}{}}
\newlabel{topologies-remark-choice-sites@cref}{{[originalreference][0][]03FF}{[1][0][]source}{}{}{}}
\makeatother
\title{Etale algebraic stacks in the common site are fppf algebraic stacks}
\author{The Stacks Project Authors}
\date{stacks-project:criteria.tex:local-snapshot-20261012; GNU FDL 1.2 or later}
\begin{document}
\maketitle
\noindent\textit{Editorial note.} No omission markers in-span; proof complete as written.\par\medskip
\noindent\textit{Allowed standard prerequisites (editorial).}\begin{itemize}
\item Groupoid presentation. algebraic-lemma-map-space-into-stack, algebraic-lemma-stack-presentation (cross-chapter, named). Canonical presentation [U/R]\_etale -> X.
\item Etale dominates smooth. more-morphisms-lemma-etale-dominates-smooth, stacks-sheaves-lemma-surjective-flat-locally-finite-presentation (cross-chapter, named). The claim that smooth covers refine etale-locally.
\item Algebraic quotient stack. algebraic-theorem-smooth-groupoid-gives-algebraic-stack (cross-chapter, named). fppf quotient stack [U/R] is algebraic.
\item ff plus es criterion. stacks-lemma-characterize-essentially-surjective-when-ff (cross-chapter, named). Concludes the equivalence.
\end{itemize}\medskip
\ifdefined\hypertarget\hypertarget{sourcepointer-1}{}\fi
\noindent\textit{Source reference.} Original reference 026O: focused-seeds/geometry/S9-crit-stacks-etale/ref-tags.txt, source lines 5--5. Navigation pointer only: the counted proof cites this statement by \textbackslash{}ref as a standing input; the deduction the counted proof has to make is carried out inside the counted span, and no part of that deduction is delegated to this pointer.\par\medskip
\ifdefined\hypertarget\hypertarget{sourcepointer-2}{}\fi
\noindent\textit{Source reference.} Original reference 04T4: focused-seeds/geometry/S9-crit-stacks-etale/ref-tags.txt, source lines 6--6. Navigation pointer only: the counted proof cites this statement by \textbackslash{}ref as a standing input; the deduction the counted proof has to make is carried out inside the counted span, and no part of that deduction is delegated to this pointer.\par\medskip
\ifdefined\hypertarget\hypertarget{sourcepointer-3}{}\fi
\noindent\textit{Source reference.} Original reference 04T5: focused-seeds/geometry/S9-crit-stacks-etale/ref-tags.txt, source lines 7--7. Navigation pointer only: the counted proof cites this statement by \textbackslash{}ref as a standing input; the deduction the counted proof has to make is carried out inside the counted span, and no part of that deduction is delegated to this pointer.\par\medskip
\ifdefined\hypertarget\hypertarget{sourcepointer-4}{}\fi
\noindent\textit{Source reference.} Original reference 04TK: focused-seeds/geometry/S9-crit-stacks-etale/ref-tags.txt, source lines 8--8. Navigation pointer only: the counted proof cites this statement by \textbackslash{}ref as a standing input; the deduction the counted proof has to make is carried out inside the counted span, and no part of that deduction is delegated to this pointer.\par\medskip
\ifdefined\hypertarget\hypertarget{sourcepointer-5}{}\fi
\noindent\textit{Source reference.} Original reference 076M: focused-seeds/geometry/S9-crit-stacks-etale/ref-tags.txt, source lines 9--9. Navigation pointer only: the counted proof cites this statement by \textbackslash{}ref as a standing input; the deduction the counted proof has to make is carried out inside the counted span, and no part of that deduction is delegated to this pointer.\par\medskip
\ifdefined\hypertarget\hypertarget{sourcepointer-6}{}\fi
\noindent\textit{Source reference.} Original reference 055V: focused-seeds/geometry/S9-crit-stacks-etale/ref-tags.txt, source lines 10--10. Navigation pointer only: the counted proof cites this statement by \textbackslash{}ref as a standing input; the deduction the counted proof has to make is carried out inside the counted span, and no part of that deduction is delegated to this pointer.\par\medskip
\ifdefined\hypertarget\hypertarget{sourcepointer-7}{}\fi
\noindent\textit{Source reference.} Original reference 025Y: focused-seeds/geometry/S9-crit-stacks-etale/ref-tags.txt, source lines 11--11. Navigation pointer only: the counted proof cites this statement by \textbackslash{}ref as a standing input; the deduction the counted proof has to make is carried out inside the counted span, and no part of that deduction is delegated to this pointer.\par\medskip
\ifdefined\hypertarget\hypertarget{sourcepointer-8}{}\fi
\noindent\textit{Source reference.} Original reference 04MB: focused-seeds/geometry/S9-crit-stacks-etale/ref-tags.txt, source lines 12--12. Navigation pointer only: the counted proof cites this statement by \textbackslash{}ref as a standing input; the deduction the counted proof has to make is carried out inside the counted span, and no part of that deduction is delegated to this pointer.\par\medskip
\ifdefined\hypertarget\hypertarget{sourcepointer-9}{}\fi
\noindent\textit{Source reference.} Original reference 046N: focused-seeds/geometry/S9-crit-stacks-etale/ref-tags.txt, source lines 13--13. Navigation pointer only: the counted proof cites this statement by \textbackslash{}ref as a standing input; the deduction the counted proof has to make is carried out inside the counted span, and no part of that deduction is delegated to this pointer.\par\medskip
\ifdefined\hypertarget\hypertarget{sourcepointer-10}{}\fi
\noindent\textit{Source reference.} Original reference 06XH: focused-seeds/geometry/S9-crit-stacks-etale/ref-tags.txt, source lines 14--14. Navigation pointer only: the counted proof cites this statement by \textbackslash{}ref as a standing input; the deduction the counted proof has to make is carried out inside the counted span, and no part of that deduction is delegated to this pointer.\par\medskip
\ifdefined\hypertarget\hypertarget{sourcepointer-11}{}\fi
\noindent\textit{Source reference.} Original reference 06TN: focused-seeds/geometry/S9-crit-stacks-etale/ref-tags.txt, source lines 15--15. Navigation pointer only: the counted proof cites this statement by \textbackslash{}ref as a standing input; the deduction the counted proof has to make is carried out inside the counted span, and no part of that deduction is delegated to this pointer.\par\medskip
\ifdefined\hypertarget\hypertarget{sourcepointer-12}{}\fi
\noindent\textit{Source reference.} Original reference 03FF: focused-seeds/geometry/S9-crit-stacks-etale/ref-tags.txt, source lines 16--16. Navigation pointer only: the counted proof cites this statement by \textbackslash{}ref as a standing input; the deduction the counted proof has to make is carried out inside the counted span, and no part of that deduction is delegated to this pointer.\par\medskip
\begin{lemma}
\label{lemma-stacks-etale}
Denote the common underlying category of $\Sch_{fppf}$
and $\Sch_\etale$ by $\Sch_\alpha$ (see
Sheaves on Stacks, Section \ref{stacks-sheaves-section-sheaves} and
Topologies, Remark \ref{topologies-remark-choice-sites}). Let $S$ be an object
of $\Sch_\alpha$. Let
$$
p : \mathcal{X} \to \Sch_\alpha/S
$$
be a category fibred in groupoids with the following properties:
\begin{enumerate}
\item $\mathcal{X}$ is a stack in groupoids over $(\Sch/S)_\etale$,
\item the diagonal $\Delta : \mathcal{X} \to \mathcal{X} \times \mathcal{X}$
is representable by algebraic spaces\footnote{Here we can either mean
sheaves in the \'etale topology whose diagonal is representable and which
have an \'etale surjective covering by a scheme or algebraic spaces as
defined in
Algebraic Spaces, Definition \ref{spaces-definition-algebraic-space}.
Namely, by Bootstrap, Lemma \ref{bootstrap-lemma-spaces-etale}
there is no difference.}, and
\item there exists $U \in \Ob(\Sch_\alpha/S)$
and a $1$-morphism $(\Sch/U)_\etale \to \mathcal{X}$
which is surjective and smooth.
\end{enumerate}
Then $\mathcal{X}$ is an algebraic stack in the sense of
Algebraic Stacks, Definition \ref{algebraic-definition-algebraic-stack}.
\end{lemma}

\begin{proof}
Note that properties (2) and (3) of the lemma and the corresponding
properties (2) and (3) of
Algebraic Stacks, Definition \ref{algebraic-definition-algebraic-stack}
are independent of the topology. This is true because these properties
involve only the notion of a $2$-fibre product of categories fibred in
groupoids, $1$- and $2$-morphisms of categories fibred in groupoids, the
notion of a $1$-morphism of categories fibred in groupoids representable
by algebraic spaces, and what it means for such a $1$-morphism to be
surjective and smooth.
Thus all we have to prove is that an \'etale stack in groupoids
$\mathcal{X}$ with properties (2) and (3) is also an fppf stack in groupoids.

\medskip\noindent
Using (2) let $R$ be an algebraic space representing
$$
(\Sch_\alpha/U) \times_\mathcal{X} (\Sch_\alpha/U)
$$
By (3) the projections $s, t : R \to U$ are smooth. Exactly as in the proof of
Algebraic Stacks, Lemma \ref{algebraic-lemma-map-space-into-stack}
there exists a groupoid in spaces $(U, R, s, t, c)$ and a canonical
fully faithful $1$-morphism $[U/R]_\etale \to \mathcal{X}$
where $[U/R]_\etale$ is the \'etale stackification of presheaf
in groupoids
$$
T \longmapsto (U(T), R(T), s(T), t(T), c(T))
$$
Claim: If $V \to T$ is a surjective smooth morphism from an algebraic space
$V$ to a scheme $T$, then there exists an \'etale covering $\{T_i \to T\}$
refining the covering $\{V \to T\}$. This follows from
More on Morphisms, Lemma \ref{more-morphisms-lemma-etale-dominates-smooth}
or the more general
Sheaves on Stacks, Lemma
\ref{stacks-sheaves-lemma-surjective-flat-locally-finite-presentation}.
Using the claim and arguing exactly as in
Algebraic Stacks, Lemma \ref{algebraic-lemma-stack-presentation}
it follows that $[U/R]_\etale \to \mathcal{X}$ is an
equivalence.

\medskip\noindent
Next, let $[U/R]$ denote the quotient stack in the fppf topology
which is an algebraic stack by
Algebraic Stacks, Theorem
\ref{algebraic-theorem-smooth-groupoid-gives-algebraic-stack}.
Thus we have $1$-morphisms
$$
U \to [U/R]_\etale \to [U/R].
$$
Both $U \to [U/R]_\etale \cong \mathcal{X}$ and
$U \to [U/R]$ are surjective and smooth (the first by assumption
and the second by the theorem) and in both cases the
fibre product $U \times_\mathcal{X} U$ and $U \times_{[U/R]} U$
is representable by $R$. Hence the $1$-morphism
$[U/R]_\etale \to [U/R]$ is fully faithful (since morphisms
in the quotient stacks are given by morphisms into $R$, see
Groupoids in Spaces, Section
\ref{spaces-groupoids-section-explicit-quotient-stacks}).

\medskip\noindent
Finally, for any scheme $T$ and morphism $t : T \to [U/R]$ the fibre product
$V = T \times_{U/R} U$ is an algebraic space surjective and smooth over $T$.
By the claim above there exists an \'etale covering $\{T_i \to T\}_{i \in I}$
and morphisms $T_i \to V$ over $T$. This proves that the object
$t$ of $[U/R]$ over $T$ comes \'etale locally from $U$. We conclude that
$[U/R]_\etale \to [U/R]$ is an equivalence of stacks in
groupoids over $(\Sch/S)_\etale$ by
Stacks, Lemma \ref{stacks-lemma-characterize-essentially-surjective-when-ff}.
This concludes the proof.
\end{proof}

\end{document}
